% !TeX root = ../main.tex

% ==================================================
% D.1. pandas_and_numpy_fundamentals
% ==================================================

\subsection{Pandas와 NumPy의 역할 구분}

\subsubsection{이번 단원의 핵심 질문}

Python에서 데이터를 다룬다는 것은 궁극적으로 객체를 생성하고, 그 객체에 연산을 적용한다는 의미다.

그런데 같은 숫자 데이터를 저장하더라도 NumPy의 \zsh{ndarray}와 Pandas의 \zsh{Series}, \zsh{DataFrame}은 서로 다른 연산 규칙을 제공한다.

이번 학습에서는 특히 다섯 가지 질문을 해결할 것이다.

\begin{dhksmry}
1. NumPy가 이미 다차원 배열을 제공하는데, Pandas가 별도로 필요한 이유는
   무엇일까?

2. ndarray와 DataFrame은 내부적으로 무엇을 데이터의 구조로 취급할까?

3. Index는 단순한 행 이름일까, 아니면 실제 연산의 의미를 결정하는 객체일까?

4. 같은 shape을 가진 두 데이터 객체에 동일한 연산을 적용해도 결과가 달라질
   수 있을까?

5. 머신러닝에서는 왜 Pandas로 데이터를 관리하다가 NumPy 배열이나 PyTorch
   Tensor로 변환할까?
\end{dhksmry}

이 질문들은 앞으로 다룰 선택, 필터링, 정렬, 결합, 데이터 누수까지 연결된다.

\subsubsection{NumPy의 기본 모델 - 위치를 가진 다차원 배열}

먼저 NumPy를 살펴보자.

NumPy의 핵심 자료 구조인 \zsh{ndarray}는 일반적으로 동일한 dtype으로 해석되는 원소들을 다차원 배열 구조로 관리하는 객체다.

대표적으로 다음 정보를 사용한다.

\begin{dhktable}
|속성	|의미|
|:---|:---|
|ndim	|배열의 차원 수|
|shape|	각 차원의 길이|
|dtype|	배열 원소의 자료형 및 저장 표현|
|strides|	각 차원을 따라 이동할 때 필요한 메모리상의 바이트 간격|
|size|	전체 원소 수|
\end{dhktable}

NumPy의 핵심은 단순히 숫자를 저장하는 데 있는 것이 아니다.

차원, 위치 dtype, 메모리 배치를 기반으로 배열 연산을 정의한다는 것이 중요하다.

\begin{dhknote}
\begin{dhksmry}[centered]
독립 실행 코드 1 - ndarray의 구조
\end{dhksmry}
\end{dhknote}

\begin{dhkcode}
import numpy as np

data = np.array(
    [
        [12, 1],
        [30, 0],
        [8, 1],
    ],
    dtype=np.int64,
)

print(type(data))
print(data.ndim)
print(data.shape)
print(data.dtype)
print(data.size)

print(data[1, 0])
\end{dhkcode}

예상 출력:

\begin{dhkzsh}
<class 'numpy.ndarray'>
2
(3, 2)
int64
6
30
\end{dhkzsh}

이 배열은 3개의 행과 2개의 열을 가진 2차원 객체다.

\begin{dhkmatrix}[delim=none,style=cards]
title :: NumPy 배열 — 위치 기반 구조
rows :: 행 0, 행 1, 행 2
cols :: 열 0, 열 1
data :: [[12,1],[green:30,0],[8,1]]
explain :: data[1, 0]은 첫 번째 차원의 위치 1과 두 번째 차원의 위치 0을 선택하므로 30을 반환한다.
\end{dhkmatrix}

여기서 주목할 것은 \zsh{data[1, 0]}이 데이터의 의미를 전혀 알 필요가 없다는 점이다.

첫 번째 열이 토큰 길이이고 두 번째 열이 클래스 레이블이라는 정보는 배열 자체의 위치 기반 연산 규칙에는 포함되지 않는다.

물론 NumPy에도 구조화 배열(Structured Array)처럼 필드 이름과 서로 다른 자료형을 다룰 수 있는 기능이 있다.

따라서 NumPy가 절대 컬럼 이름을 가질 수 없다고 이해해서는 안 된다.

다만 일반적인 \zsh{ndarray}에는 Pandas가 제공하는 행·열 레이블과 자동 정렬 모델이 내재되어 있지 않다.

\subsubsection{Pandas의 기본 모델 - 레이블을 가진 데이터 객체}

이번에는 똑같은 데이터에 의미를 부여해 보자.

숫자 12, 30, 8이 각각 문서의 토큰 길이라고 가정하고, 1과 0은 문서의 분류 레이블이라고 생각해 보자.

Pandas에서는 이 정보를 다음과 같이 구성할 수 있다.

\begin{dhkcode}
import pandas as pd

df = pd.DataFrame(
    {
        "token_length": [12, 30, 8],
        "label": [1, 0, 1],
    },
    index=["doc-A", "doc-B", "doc-C"],
)

print(df)
\end{dhkcode}

예상 출력:

\begin{dhkzsh}
       token_length  label
doc-A            12      1
doc-B            30      0
doc-C             8      1
\end{dhkzsh}

NumPy 배열과 비교하면 다음 정보가 추가되었다.

\begin{dhksmry}
   - 행마다 doc-A, doc-B, doc-C라는 레이블이 있다.

   - 열마다 token_length, label이라는 레이블이 있다.

   - 이 레이블을 이용하여 데이터를 조회하거나 연산할 수 있다.
\end{dhksmry}

하지만 이러한 레이블이 단순히 출력용 문자로 붙어 있는 것은 아니다.

Pandas는 레이블을 데이터 구조의 일부로 관리하고, 여러 연산에서 레이블 사이의 대응 관계를 사용한다.

\begin{dhknote}
\begin{dhksmry}[centered]
객체 구조
\end{dhksmry}
\end{dhknote}

\begin{dhkgraph}[uniform]
node n1 (1,2) :: DataFrame 객체 // 2차원 레이블 데이터 구조
node n2 (2,1) :: index // 행 레이블 // doc-A // doc-B // doc-C
node n3 (2,2) :: columns // 열 레이블 // token_length // label
node n4 (2,3) :: 데이터 // 열별 값과 dtype // 12, 30, 8 // 1, 0, 1
node n5 (3,2) :: 레이블 + 타입화된 데이터 // 선택·정렬·계산·결합에 사용

split n1 -> n2, n3, n4
edge n3 -- n5

caption :: 개념적 객체 구조 // 실제 Pandas의 메모리 저장 방식이나 내부 관리 객체를 그대로 묘사한 그림은 아니다.
\end{dhkgraph}

여기에서 기존 Python 객체 모델과 연결할 중요한 점이 있다.

\zsh{df}라는 이름이 DataFrame을 담고 있는 것은 아니다.

\zsh{df}는 생성된 \zsh{DataFrame} 인스턴스를 참조하는 이름이다.

그리고 \zsh{df.index}와 \zsh{df.columns} 역시 독립적인 데이터 구조를 나타내는 속성이며, 일반적인 Python 리스트로 단순화할 수 없다.

이 관계는 \zsh{D.2}에서 구체적으로 분석할 것이다.

\subsubsection{ndarray, Series, DataFrame, Index의 관계}

네 객체는 모두 데이터 처리를 지원하지만 구조적 역할이 다르다.

\begin{dhktable}
|객체|차원|레이블|dtype 구조|
|:---|:---|:---|:---|
|ndarray|0차원 이상|기본적으로 위치 중심|일반적으로 배열 전체에 하나의 dtype|
|Series|1차원|하나의 Index|Series 하나에 하나의 논리적 dtype|
|DataFrame|2차원|행 Index, 열 Index|열마다 다른 dtype 가능|
|Index|일반적으로 1차원 레이블 구조|레이블 자체를 관리|Index의 자료형을 가짐|
\end{dhktable}

여기서 \zsh{DataFrame}을 단순히 2차원 배열이라고 부르는 것은 충분하지 않다.

\zsh{DataFrame}은 열마다 다른 dtype을 가질 수 있고, 각 축의 레이블을 유지하면서 연산을 수행할 수 있기 때문이다.

또한 DataFrame을 개념적으로 여러 Series의 집합처럼 이해할 수는 있지만, 실제 모든 내부 데이터가 독립적인 NumPy 배열이나 Series 객체로 저장된다고 단정해서는 안 된다.

Pandas는 ExtensionArray 등 다양한 저장 표현을 지원한다.

\begin{dhknote}
\begin{dhksmry}[centered]
독립 실행 코드 2 - 객체와 dtype 확인
\end{dhksmry}
\end{dhknote}

\begin{dhkcode}
import numpy as np
import pandas as pd

np_data = np.array(
    [[12, 1], [30, 0], [8, 1]],
    dtype=np.int64,
)

df = pd.DataFrame(
    {
        "token_length": [12, 30, 8],
        "label": [1, 0, 1],
        "language": ["ko", "en", "ko"],
    },
    index=pd.Index(
        ["doc-A", "doc-B", "doc-C"],
        name="doc_id",
    ),
).astype({
    "token_length": "int64",
    "label": "int64",
    "language": "object",
})

print(type(np_data).__name__)
print(np_data.shape, np_data.ndim, np_data.dtype)

print(type(df).__name__)
print(df.shape, df.ndim)

print(df.index.tolist())
print(df.columns.tolist())

print({
    name: str(dtype)
    for name, dtype in df.dtypes.items()
})

print(type(df["token_length"]).__name__)
print(df["token_length"].index.tolist())
\end{dhkcode}

\begin{dhknote}
\begin{dhksmry}[centered]
실행 과정 분석
\end{dhksmry}
\end{dhknote}

첫 번째: \zsh{np.array(...)}

\zsh{np.array}는 NumPy 모듈에 정의된 함수다.

호출 결과로 \zsh{ndarray} 인스턴스를 반환하며, 그 객체에 \zsh{np_data}라는 이름이 바인딩된다.

배열 전체에 \zsh{int64} dtype을 지정했기 때문에 두 열 모두 같은 dtype으로 저장된다.

두 번째: \zsh{pd.DataFrame(...)}

\zsh{pd.DataFrame}은 클래스 객체다.

이 클래스 객체를 호출하면 전달된 데이터와 인덱스를 바탕으로 새로운 DataFrame 인스턴스가 생성된다.

이어지는 \zsh{.astype(...)}은 변환된 DataFrame을 반환하고, 최종 결과가 \zsh{df}에 바인딩된다.

이 코드에서는 \zsh{pd.DataFrame(...)} 호출의 중간 결과를 별도의 이름으로 보존하지 않았다.

세 번째: \zsh{df.dtypes}

NumPy 배열에서는 보통 \zsh{data.dtype}으로 배열 원소의 dtype을 확인한다.

반면 DataFrame에서는 각 열의 dtype을 표시하는 \zsh{df.dtypes}를 사용할 수 있다.

그 반환값은 열 레이블을 Index로 사용하는 Series다.

네 번째: \zsh{df["token_length"]}

이 표현식에서 \zsh{[]}는 일반적인 Python 인덱싱 문법이지만, 실제 선택 규칙은 DataFrame이 제공하는 \zsh{__getitem__()} 구현에 의해 결정된다.

여기서는 \zsh{token_length}라는 열 레이블에 대응하는 Series를 반환한다.

그 Series는 해당 열의 값을 다루면서 행 Index도 유지한다.

중요한 점은 반환된 Series와 원본 DataFrame이 내부 데이터를 공유하는지, 이후 수정이 원본에 영향을 미치는지는 별도의 문제라는 것이다.

이 부분은 \zsh{D.7}에서 Copy-on-Write와 연결해 분석할 것이다.

\subsubsection{NumPy의 axis와 Pandas의 axis}

이제 NumPy 선행 개념에서 가장 중요한 \zsh{axis}를 연결해 보자.

NumPy에서 \zsh{axis}는 배열의 차원을 가리킨다.

2차원 배열이라면 일반적으로 다음과 같이 해석한다.

\begin{dhksmry}
   - axis=0 : 첫 번째 차원, 보통 행 위치가 나열되는 축
   - axis=1 : 두 번째 차원, 보통 열 위치가 나열되는 축
\end{dhksmry}

하지만 집계 연산에서 중요한 것은

\begin{dhksmry}[centered]
지정한 축을 따라 원소를 집계한다
\end{dhksmry}

는 점이다.

따라서 \zsh{sum(axis=0)}은 첫 번째 차원을 따라 합산해 해당 차원을 축소하고,

\zsh{sum(axis=1)}은 두 번째 차원을 따라 합산한다.

Pandas에서도 2차원 DataFrame의 \zsh{axis=0, axis=1}은 이 차원 구조와 대응하지만, 축에 레이블이 부여되어 있다.

\zsh{axis="index"}와 \zsh{axis="columns"}라는 표현도 사용할 수 있다.

\begin{dhknote}
\begin{dhksmry}[centered]
독립 실행 코드 3 - 축에 따른 집계
\end{dhksmry}
\end{dhknote}

\begin{dhkcode}
import numpy as np
import pandas as pd

arr = np.array(
    [[2, 1], [4, 0]],
    dtype=np.int64,
)

df = pd.DataFrame(
    arr,
    index=["doc-A", "doc-B"],
    columns=["tokens", "label"],
)

print(arr.sum(axis=0).tolist())
print(df.sum(axis="index").to_dict())

print(arr.sum(axis=1).tolist())
print(df.sum(axis="columns").to_dict())
\end{dhkcode}

예상 출력:

\begin{dhkzsh}
[6, 1]
{'tokens': 6, 'label': 1}
[3, 4]
{'doc-A': 3, 'doc-B': 4}
\end{dhkzsh}

여기에서 조심할 부분이 있다.

\zsh{axis=0}을 무조건 "행을 선택한다"라고 외우면 안 된다.

집계에서는 행을 따라 값을 모으지만, 다른 API에서는 지정한 축에 레이블을 적용하거나 그 축을 따라 연산하도록 하는 등 구체적인 의미가 달라질 수 있다.

따라서 axis가 어느 차원을 가리키는지와 해당 메서드가 그 축을 어떻게 사용하는지를 함께 확인해야 한다.

\subsubsection{가장 중요한 차이 - 위치 연산과 Index alignment}

이제 Pandas를 배우는 이유가 명확해지는 예제를 살펴보자.

NumPy 배열 두 개를 더한다고 생각해 보자.

\begin{dhkcode}
import numpy as np

a = np.array([10, 20])
b = np.array([1, 2])

print((a + b).tolist())
\end{dhkcode}

예상 출력:

\begin{dhkzsh}
[11, 22]
\end{dhkzsh}

브로드캐스팅이 필요 없는 이 예제에서는 같은 위치의 원소끼리 더한다.

그렇다면 똑같은 값으로 Pandas Series를 만들면 어떻게 될까?

\begin{dhknote}
\begin{dhksmry}[centered]
독립 실행 코드 4 - Series의 정렬 연산
\end{dhksmry}
\end{dhknote}

\begin{dhkcode}
import numpy as np
import pandas as pd

a = np.array([10, 20], dtype=np.int64)
b = np.array([1, 2], dtype=np.int64)

left = pd.Series(
    a,
    index=["doc-A", "doc-B"],
    name="score",
)

right = pd.Series(
    b,
    index=["doc-B", "doc-A"],
    name="bonus",
)

result = left + right

print((a + b).tolist())
print(result.to_dict())
print(left.to_dict())
print(result is left)
\end{dhkcode}

예상 출력:

\begin{dhkzsh}
[11, 22]
{'doc-A': 12, 'doc-B': 21}
{'doc-A': 10, 'doc-B': 20}
False
\end{dhkzsh}

왜 Pandas에서는 \zsh{[11, 22]}가 아니라 \zsh{doc-A}에 12, \zsh{doc-B}에 21이 대응할까?

두 Series의 값을 정확히 살펴보자.

\begin{dhktable}
|Index|	left|	right|	결과|
|:---|:---|:---|:---|
|doc-A|	10|	2|	12|
|doc-B|	20|	1|	21|
\end{dhktable}

\zsh{right}의 첫 번째 값은 1이지만, 그 값에 대응하는 레이블은 \zsh{doc-B}다.

따라서 Pandas는 그 값을 \zsh{left}의 첫 번째 위치가 아니라 \zsh{doc-B}에 대응시켜 계산한다.

\begin{dhknote}
\begin{dhksmry}[centered]
Python 객체 모델로 실행 추적하기
\end{dhksmry}
\end{dhknote}

이 연산을 기존 Python의 방식으로 분석해 보자.

\begin{dhknote}
\begin{dhksmry}[centered]
1단계 - 객체 생성
\end{dhksmry}
\end{dhknote}

\zsh{pd.Series(...)} 호출을 통해 서로 다른 Series 객체 두 개가 생성된다.

각 객체는 값뿐 아니라 Index도 가지고 있다.

\begin{dhknote}
\begin{dhksmry}[centered]
2단계 - 이름 바인딩
\end{dhksmry}
\end{dhknote}

\zsh{left}와 \zsh{right}은 각각 생성된 Series 객체에 바인딩된다.

\begin{dhknote}
\begin{dhksmry}[centered]
3단계 - 이항 연산
\end{dhksmry}
\end{dhknote}

\zsh{left + right}에서 \zsh{+} 연산자는 Python의 특수 메서드 기반 연산 프로토콜을 사용한다.

Pandas의 Series 연산 구현은 레이블 정렬 규칙을 적용하고, 대응하는 값을 계산한다.

따라서 NumPy 배열의 위치 기반 덧셈과 다른 결과가 발생한다.

\begin{dhknote}
\begin{dhksmry}[centered]
4단계 - 반환 객체
\end{dhksmry}
\end{dhknote}

이 예제의 덧셈 결과는 새로운 Series 객체다.

\begin{dhkcode}
result is left
\end{dhkcode}

가 \zsh{False}인 이유는 \zsh{result}가 \zsh{left}와 다른 객체를 참조하기 때문이다.

\begin{dhknote}
\begin{dhksmry}[centered]
5단계 - 상태 변화
\end{dhksmry}
\end{dhknote}

\zsh{left + right} 자체는 \zsh{left}의 기존 값을 변경하지 않는다.

확인 결과 \zsh{left}는 여전히 다음 데이터를 가지고 있다.

\begin{dhkzsh}
{'doc-A': 10, 'doc-B': 20}
\end{dhkzsh}

여기서 명심할 점이 있다.

Series 간 산술 연산에서 새 객체를 반환한다는 사실과 모든 Pandas 연산이 원본을 변경하지 않는다는 주장은 전혀 다르다.

명시적인 대입, 변경 메서드, Copy-on-Write 정책은 별도로 분석해야 한다.

이 주제는 \zsh{D.7}에서 다룰 것이다.

\begin{dhknote}
\begin{dhksmry}[centered]
예외와 경계 조건
\end{dhksmry}
\end{dhknote}

Pandas가 언제나 모든 데이터를 행 레이블로 정렬하는 것은 아니다.

예를 들어 \zsh{DataFrame + Series}에서는 기본적으로 Series의 Index를

DataFrame의 열 레이블과 정렬해 연산한다.

DataFrame끼리의 연산에서는 행과 열 레이블을 모두 고려할 수 있다.

따라서 정확한 판단 기준은 다음과 같다.

\begin{dhksmry}[centered]
어떤 객체 타입들이 연산에 참여하고,

그 연산에서 어느 축의 레이블을 대응시키는가?
\end{dhksmry}

이 질문을 \zsh{D.5}에서 심화 학습 할 것이다.

\subsubsection{머신러닝에서 NumPy와 Pandas를 함께 사용하는 이유}

지금까지의 차이를 실제 Transformer 데이터 파이프라인과 연결해 보자.

Transformer 모델에는 일반적으로 정수 토큰 ID 배열이 필요하다.

예를 들어 두 문서를 토큰화했다고 가정하자.

\begin{dhkcode}
import numpy as np

token_ids = np.array(
    [
        [101, 202, 0],
        [101, 303, 404],
    ],
    dtype=np.int64,
)
\end{dhkcode}

여기서

\begin{dhksmry}
   - 첫 번째 차원은 샘플

   - 두 번째 차원은 토큰 위치
\end{dhksmry}

를 나타낸다.

0은 예제에서 패딩 토큰 ID로 사용했다고 가정한다.

그런데 학습을 진행하려면 토큰 ID 이외의 정보도 필요하다.

문서 식별자, 정답 레이블, 원본 언어, 실제 토큰 길이 등이 대표적이다.

이 정보를 관리하기 위해 Pandas를 사용할 수 있다.

\begin{dhknote}
\begin{dhksmry}[centered]
독립 실행 코드 5 - 모델 입력과 메타데이터
\end{dhksmry}
\end{dhknote}

\begin{dhkcode}
import numpy as np
import pandas as pd

# 모델 입력: 문서 순서 doc-B, doc-A
token_ids = np.array(
    [
        [101, 202, 0],
        [101, 303, 404],
    ],
    dtype=np.int64,
)

# 분석용 메타데이터
metadata = pd.DataFrame(
    {
        "label": [0, 1],
        "token_length": [2, 3],
        "language": ["ko", "en"],
    },
    index=pd.Index(
        ["doc-B", "doc-A"],
        name="doc_id",
    ),
)

labels = metadata["label"].to_numpy(
    dtype=np.int64
)

print(token_ids.shape)
print(labels.shape)

print(metadata.index.tolist())
print(labels.tolist())
\end{dhkcode}

예상 출력:

\begin{dhkzsh}
(2, 3)
(2,)
['doc-B', 'doc-A']
[0, 1]
\end{dhkzsh}

이 예제에서는 토큰 배열의 0번째 행과 메타데이터의 첫 번째 행이 모두 \zsh{doc-B}라는 규약을 사용하고 있다.

현재까지는 문제가 없다.

그런데 다음 코드를 실행하면 어떻게 될까?

\begin{dhkcode}
sorted_metadata = metadata.sort_index()

print(sorted_metadata.index.tolist())
print(
    sorted_metadata["label"]
    .to_numpy(dtype=np.int64)
    .tolist()
)

print(
    metadata.index.equals(
        sorted_metadata.index
    )
)
\end{dhkcode}

예상 출력:

\begin{dhkzsh}
['doc-A', 'doc-B']
[1, 0]
False
\end{dhkzsh}

여기서 매우 중요한 문제가 발생했다.

\zsh{metadata}는 정렬되었지만 \zsh{token_ids}는 정렬되지 않았다.

따라서 정렬된 DataFrame에서 레이블을 추출해 원래 토큰 배열에 그대로 연결한다면 잘못된 정답이 대응할 수 있다.

그런데도 두 객체의 shape은 각각 \zsh{(2, 3)}과 \zsh{(2,)}이므로 일반적인 모델 학습 코드에서는 형태 오류가 발생하지 않을 수 있다.

이 문제를 더 엄밀하게 표현하면 다음과 같다.

\[
\text{Shape Compatibility}
\;\not\Rightarrow\;
\text{Sample Identity Consistency}
\]

즉, 배열의 크기가 호환된다고 해서 샘플 식별자의 대응 관계까지 보장되는 것은 아니다.

이것이 Pandas의 레이블 관리가 실제 머신러닝에서 중요한 이유 중 하나다.

실전에서는 정렬·선택·결합 이후에도 공통 샘플 ID와 행 순서를 검증해야 한다.

이 부분은 \zsh{D.13}의 결합과 \zsh{D.21}의 모델 입력 변환에서 다시 다룰 것이다.

\subsubsection{특히 주의해야 할 세 가지 오해}

\begin{dhknote}
\begin{dhksmry}[centered]
첫째, Pandas는 NumPy에 행과 열 이름만 붙인 라이브러리다?
\end{dhksmry}
\end{dhknote}

그렇게 설명하면 불완전하다.

Pandas는 레이블 정렬, 열별 dtype, 결측 표현, 그룹화, 데이터 결합 등 별도의 데이터 처리 의미론을 제공한다.

NumPy의 일부 기능과 저장 구조를 활용하지만, 단순히 \zsh{ndarray}를 보기 좋게 표시하는 도구는 아니다.

\begin{dhknote}
\begin{dhksmry}[centered]
둘째, Index가 존재하면 행을 유일하게 식별할 수 있다?
\end{dhksmry}
\end{dhknote}

반드시 그렇지는 않다.

Pandas는 중복된 Index 레이블도 허용한다.

\begin{dhkcode}
import pandas as pd

df = pd.DataFrame(
    {"label": [0, 1]},
    index=["doc-A", "doc-A"],
)

print(df.index.is_unique)
print(type(df.loc["doc-A"]).__name__)
\end{dhkcode}

예상 출력:

\begin{dhkzsh}
False
DataFrame
\end{dhkzsh}

동일한 레이블을 가진 행이 두 개 있기 때문에 \zsh{df.loc["doc-A"]}는 하나의 행만 반환하지 않는다.

따라서 Index는 레이블 체계이지, 자동으로 보장되는 기본키(Primary Key)는 아니다.

어떤 연산은 중복 Index를 허용하지만, 유일성을 요구하는 연산은 오류를 발생시킬 수 있다.

\begin{dhknote}
\begin{dhksmry}[centered]
셋째, DataFrame을 NumPy 배열로 변환하면 
같은 데이터이므로 의미도 보존된다?
\end{dhksmry}
\end{dhknote}

수치 값은 보존할 수 있지만 레이블 정보는 일반적으로 NumPy 배열의 축 레이블로 전달되지 않는다.

또한 여러 dtype으ㅢ 열을 하나의 NumPy 배열로 결합하는 과정에서 공통 dtype으로 변환되거나 \zsh{object} dtype이 사용될 수 있다.

그리고 \zsh{.to_numpy()}가 항상 메모리 복사를 수행하거나, 항상 원본 메모리를 공유한다고 가정해서도 안 된다.

이것들은 각각 값의 변환, dtype의 변환, 레이블 정보의 소실, 메모리 공유라는 별개의 문제다.

\subsubsection{Pandas 버전별 동작의 경계}

Pandas 3.0.0이 출시되면서 두 가지 큰 변화가 있었다.

\begin{dhktable}
|항목|	Pandas 2.x의 일반적 기본 동작|	Pandas 3.0 이상|
|:---|:---|:---|
|문자열 dtype 추론|	일반적으로 object|	기본적으로 전용 str dtype|
|Copy-on-Write|	버전과 설정에 의존|	기본이자 유일한 복사·뷰 동작 모드|
|연쇄 대입|	경우에 따라 경고 및 불확실한 결과|	원본 변경 방식으로 사용할 수 없음|
\end{dhktable}

앞에서 작성한 코드에 \zsh{astype({"language": "object"})}를 명시한 이유도 문자열 dtype 추론에 따른 버전 차이를 제거하기 위해서다.

이번 \zsh{D.1}에서는 그 변화가 존재한다는 사실까지만 기억하면 충분하다.

\subsubsection{핵심 정리와 판단 기준}

지금까지 배운 내용을 비교하면 다음과 같다.

\begin{dhktable}
|질문|	NumPy|	Pandas|
|:---|:---|:---|
|핵심 추상화는?|	다차원 배열|	레이블을 가진 데이터 객체|
|기본적인 접근 기준은?|	차원과 위치|	위치와 레이블|
|이항 연산의 중요한 규칙은?|	shape과 브로드캐스팅|	객체 타입에 따른 레이블 정렬 및 브로드캐스팅|
|자료형은 어떻게 관리하나?|	일반적으로 배열 단위|	Series 단위, DataFrame에서는 열 단위|
|행 식별자를 자동 관리하나?|	별도 관리 필요|	Index 제공, 단 유일성 보장은 별도|
|ML/DL에서 대표적 용도는?|	수치 배열과 계산용 데이터|	데이터 정제·메타데이터·집계·정합성 관리|
\end{dhktable}

이번 학습에서 도달해야 하는 핵심 결론은 다음과 같다.

\begin{dhksmry}[centered]
NumPy는 배열의 차원·위치·dtype을 중심으로 수치 연산을 구성하고,

Pandas는 여기에 데이터 레이블과 열별 자료형 등을 제공한다.
\end{dhksmry}

그리고 가장 중요한 것은 다음 판단 원칙이다.

\begin{dhksmry}[centered]
두 데이터 객체의 shape이 같다는 사실만으로 두 객체의 값들이 

동일한 의미의 관측치에 대응한다고 결론 내릴 수는 없다.
\end{dhksmry}

이 원칙은 앞으로 배우게 될 Index alignment, JOIN, GroupBy, 데이터 누수, PyTorch Dataset 생성까지 계속 적용된다.